%! Author = sriadityavedantam
%! Date = 3/29/26

\section{Units and Divisors}

\begin{definition}[Units in Congruence Classes]
    $[a]$ is a unit if $[a]$ has a multiplicative inverse.
    That is, there exists $[x]\in \zn$ such that
    \[
        [a][x] = [1].
    \]
\end{definition}

\begin{theorem}{
    $[a]$ is a unit if and only if $\gcd(a,n) = 1$.
}
\end{theorem}

\begin{proposition}{
    All nonzero classes in $\zp$ are units.
}
    Suppose $[a]\in\zp$ with $[a]\neq [0]$.
    Then $p\nmid a$.
    Since $p$ is prime, this implies
    \[
        \gcd(a,p) = 1.
    \]
    By Bézout's Identity, there exist integers $x,q$ such that
    \[
        ax + qp = 1.
    \]
    Thus
    \[
        ax \equiv 1 \mod p,
    \]
    so
    \[
        [a][x] = [1].
    \]
    Therefore $[a]$ is a unit.
\end{proposition}

Easy to show the converse by showing a multiplicative inverse in $\zp$.
So $[a]$ has a multiplicative inverse.
This proposition, using $n\neq p$, will show it is true for the previous theorem.
To test whether $[a]$ is a unit in $\mathbb{Z}_{32}$ and find $[a]^{-1}$ in $\mathbb{Z}_{32}$, let $a=4$ and try to find an $x\in\z$ such that
\[
    4x + 32q = 1.
\]
But
\[
    \gcd(4,32) = 4 \neq 1,
\]
so no such integers $x,q$ exist.
Therefore $[4]$ is not a unit in $\mathbb{Z}_{32}$.
So there is no inverse to find.

\begin{definition}[Zero-Divisors]
    $[a]$ is a zero-divisor in $\zn$ if there exists $[x]\neq [0]$ such that
    \[
        [ax] = [0].
    \]
\end{definition}

\begin{theorem}{
    $[a]$ is a zero-divisor in $\zn$ if and only if $\gcd(a,n)\neq 1$.
}
Let us prove the contrapositive in one direction.
Suppose $\gcd(a,n) = 1$.
Then $[a]$ is a unit in $\zn$ by the previous theorem.
We will show that $[a]$ is not a zero-divisor.
Suppose $[ab] = [0]$ for some $b\in\z$.
Since $[a]$ is a unit, there exists $[x]\in\zn$ such that
\[
    [xa] = [1].
\]
Then
\[
    [x]([ab]) = [x][0] = [0],
\]
so
\[
    [xa][b] = [1][b] = [b] = [0].
\]
Thus $[ab]=[0]$ implies $[b]=[0]$.
Therefore $[a]$ is not a zero-divisor.
Conversely, suppose $\gcd(a,n) = d > 1$.
Write
\[
    a = da_1, \qquad n = dn_1.
\]
Then
\[
    a n_1 = da_1 n_1 = a_1 n,
\]
so
\[
    [a][n_1] = [0].
\]
Also, since $d>1$, we have $0<n_1<n$, hence
\[
    [n_1]\neq [0].
\]
Therefore $[a]$ is a zero-divisor in $\zn$.
\end{theorem}

For example, if we take $\mathbb{Z}_{12}$, then since
\[
    \gcd(4,12) = 4,
\]
$[4]$ is a zero-divisor.